\documentclass[aspectratio=169]{beamer}
\usepackage[utf8]{inputenc}
\usepackage[spanish]{babel}
\usepackage{listings}

% Estilo personalizado Beamer (ICESI)
\usepackage{custombeamer}

% ==============================================================================
% CONFIGURACIÓN DE LISTINGS: PALETA VIBRANTE MULTICOLOR DE ALTO CONTRASTE
% ==============================================================================
\definecolor{codeBg}{HTML}{13141F}        % Fondo grafito oscuro profundo
\definecolor{codeFrame}{HTML}{2A2D3D}     % Borde sutil del marco
\definecolor{codeText}{HTML}{FFFFFF}      % Texto plano en blanco puro
\definecolor{codeKeyword}{HTML}{FF2A85}   % Palabras reservadas en Rosa Neón (#FF2A85)
\definecolor{codePrimitive}{HTML}{00F5A0} % Tipos primitivos en Verde Menta Neón (#00F5A0)
\definecolor{codeType}{HTML}{00D2FF}      % Entidades y Clases en Cyan / Sky Blue (#00D2FF)
\definecolor{codeAnnotation}{HTML}{FFD600}% Anotaciones Spring en Amarillo / Dorado Neón (#FFD600)
\definecolor{codeBoolean}{HTML}{FF6B6B}   % Booleanos y constantes en Coral / Salmón (#FF6B6B)
\definecolor{codeString}{HTML}{A7F432}    % Cadenas de texto en Lima Neón (#A7F432)
\definecolor{codeComment}{HTML}{8E9AB0}   % Comentarios en Gris Lavanda (#8E9AB0)

\lstdefinestyle{SpringCodeStyle}{
  language=Java,
  basicstyle=\ttfamily\fontsize{6pt}{7.5pt}\selectfont\color{codeText},
  deletekeywords=[1]{boolean,int,long,void,double,float,char,byte,true,false,null},
  keywordstyle=[1]\color{codeKeyword}\bfseries,
  keywordstyle=[2]\color{codePrimitive}\bfseries,
  keywordstyle=[3]\color{codeType}\bfseries,
  keywordstyle=[4]\color{codeAnnotation}\bfseries,
  keywordstyle=[5]\color{codeBoolean}\bfseries,
  stringstyle=\color{codeString},
  commentstyle=\color{codeComment}\itshape,
  alsoletter={@},
  morekeywords=[1]{public,private,final,class,interface,return,new,extends,implements,throw,if,var},
  morekeywords=[2]{boolean,int,long,void,double,float,char,byte},
  morekeywords=[3]{Usuario,Profesor,Curso,Estudiante,EstudianteCurso,Rol,Permiso,Optional,List,Page,Pageable,String,Long,Integer,ResponseEntity,PageRequest,Sort},
  morekeywords=[4]{@Service,@Repository,@Transactional,@Query,@Param,@RequiredArgsConstructor,@Modifying,@GetMapping,@RequestParam,@PathVariable},
  morekeywords=[5]{true,false,null},
  numbers=none,
  breaklines=true,
  backgroundcolor=\color{codeBg},
  frame=single,
  rulecolor=\color{codeFrame},
  tabsize=2,
  aboveskip=1pt,
  belowskip=1pt,
  showstringspaces=false
}
\lstset{style=SpringCodeStyle}

\begin{document}

% ==============================================================================
% PORTADA & AGENDA
% ==============================================================================

\titleSlideMinimalist
  {Spring Data JPA}
  {Services, Query Methods y Consultas Avanzadas}
  {Prof. Alejandro Pe\~naranda}
  {Universidad ICESI -- Compunet II}

\agendaSlideNumbered
  {Agenda de la Sesi\'on}
  {Capa de Servicios y Repositorios}
  {Query Methods por Convenci\'on}
  {Consultas Avanzadas con @Query}
  {Estudio Aut\'onomo: Transacciones y Paginaci\'on}

% ==============================================================================
% SECCIÓN 1: CAPA DE SERVICIOS Y CONSUMO DE REPOSITORIOS
% ==============================================================================

\sectionSlideStripe{01}{Capa de Servicios}{Invocaci\'on de Repositorios y Separaci\'on de Capas}

\contentSlideBulletCard{Arquitectura en Capas: El Rol del Service}
  {
    \begin{itemize}
      \item \textbf{Aislamiento de L\'ogica:} El Controller jam\'as accede directamente a la base de datos.
      \item \textbf{Orquestaci\'on:} El Service coordina m\'ultiples repositorios y entidades.
      \item \textbf{Reutilizaci\'on:} Reglas consumibles por REST, procesos batch o eventos.
    \end{itemize}
  }
  {Principio SRP}
  {El Service encapsula las reglas de negocio y delega las operaciones de acceso a datos al Repositorio.}

\contentSlideTwoCols{Flujo de Consulta: Controller vs Service}
  {Controller (Presentaci\'on)}
  {
    \begin{itemize}
      \item Mapea rutas y verbos HTTP.
      \item Recibe \texttt{@RequestParam} / \texttt{@PathVariable}.
      \item Valida tipos de datos de entrada.
      \item Retorna \texttt{ResponseEntity<T>}.
    \end{itemize}
  }
  {Service (Negocio)}
  {
    \begin{itemize}
      \item Aplica las pol\'iticas del dominio.
      \item Invoca Query Methods del repositorio.
      \item Transforma entidades a DTOs.
      \item Orquesta m\'ultiples entidades.
    \end{itemize}
  }

\contentSlideBulletCard{Inyecci\'on de Dependencias Robusta}
  {
    \begin{itemize}
      \item \textbf{Inyecci\'on por Constructor:} Garantiza inmutabilidad con atributos \texttt{final}.
      \item \textbf{Lombok @RequiredArgsConstructor:} Genera el constructor de dependencias autom\'aticamente.
      \item \textbf{Mocks y Tests:} Permite instanciar el Service en tests aislando los Repositorios con Mockito.
    \end{itemize}
  }
  {Inmutabilidad}
  {Evitar \texttt{@Autowired} sobre campos para eliminar referencias nulas y facilitar pruebas unitarias.}

\begin{frame}[fragile]{Ejemplo: Service Consumiendo Repositorios}
\begin{lstlisting}[style=SpringCodeStyle]
@Service
@RequiredArgsConstructor
public class CursoService {
    private final CursoRepository cursoRepository;

    // Invoca un metodo basico de busqueda en el repositorio
    public List<Curso> listarCursosPorDepartamento(String depto) {
        return cursoRepository.findByDepartamento(depto);
    }

    // Invoca la busqueda por identificador unico
    public Optional<Curso> obtenerCursoPorId(Long id) {
        return cursoRepository.findById(id);
    }
}
\end{lstlisting}
\end{frame}

% ==============================================================================
% SECCIÓN 2: QUERY METHODS POR CONVENCIÓN (PASO A PASO DIDÁCTICO)
% ==============================================================================

\sectionSlideStripe{02}{Query Methods}{Introspecci\'on y Ejercicios Granulares}

\contentSlideBulletCard{?`Qu\'e es un Query Method?}
  {
    \begin{itemize}
      \item \textbf{Definici\'on:} Consulta derivada declarada \'unicamente mediante la firma de un m\'etodo en un \texttt{Repository}.
      \item \textbf{Convenci\'on de Nombres:} El nombre del m\'etodo expresa directamente el criterio de b\'usqueda sin escribir SQL.
      \item \textbf{Cero Boilerplate:} Elimina la necesidad de escribir implementaciones repetitivas de persistencia.
    \end{itemize}
  }
  {Idea Clave}
  {El nombre del m\'etodo Java se convierte en la propia sentencia de persistencia ejecutada contra la base de datos.}

\contentSlideBulletCard{?`C\'omo Funcionan los Query Methods?}
  {
    \begin{itemize}
      \item \textbf{Introspecci\'on en Arranque:} Spring analiza el nombre de cada m\'etodo de la interfaz.
      \item \textbf{Traducci\'on Autom\'atica:} Convierte nombres Java en cl\'ausulas SQL/JPQL.
      \item \textbf{Cero Implementaci\'on:} T\'u s\'olo declaras la firma; Spring genera el proxy.
    \end{itemize}
  }
  {Sin SQL Manual}
  {El desarrollador no escribe consultas manuales para el 80\% de las b\'usquedas del dominio.}

% ------------------------------------------------------------------------------
% PASO 1: BÚSQUEDA ÚNICA (findBy + Optional$<$T$>$)
% ------------------------------------------------------------------------------
\contentSlideBulletCard{B\'usqueda \'Unica: findBy... con Optional$<$T$>$}
  {
    \begin{itemize}
      \item \textbf{Prop\'osito:} Consultar un registro \'unico mediante un campo identificador o \'unico.
      \item \textbf{Retorno Optional$<$T$>$:} Modela expl\'icitamente la presencia o ausencia de la entidad.
      \item \textbf{Evita NullPointer:} El Service decide si desempaquetar o lanzar una excepci\'on controlada.
    \end{itemize}
  }
  {Claves \'Unicas}
  {Ideal para buscar por correo institucional, c\'edula o nombres \'unicos en el sistema.}

\begin{frame}[fragile]{Ejemplo 1: B\'usqueda \'Unica con Optional$<$T$>$}
\begin{lstlisting}[style=SpringCodeStyle]
// 1. Declaracion en UsuarioRepository
@Repository
public interface UsuarioRepository extends JpaRepository<Usuario, Long> {
    Optional<Usuario> findByCorreoInstitucional(String correo);
}

// 2. Invocacion limpia desde UsuarioService
@Service
@RequiredArgsConstructor
public class UsuarioService {
    private final UsuarioRepository usuarioRepository;

    public Optional<Usuario> obtenerPorCorreo(String correo) {
        return usuarioRepository.findByCorreoInstitucional(correo);
    }
}
\end{lstlisting}
\end{frame}

% ------------------------------------------------------------------------------
% PASO 2: EXISTENCIA BOOLEANA (existsBy + boolean)
% ------------------------------------------------------------------------------
\contentSlideBulletCard{Comprobaci\'on de Existencia: existsBy... con boolean}
  {
    \begin{itemize}
      \item \textbf{Prop\'osito:} Verificar r\'apidamente si ya existe un registro que cumpla el criterio.
      \item \textbf{Retorno boolean:} Devuelve \texttt{true} si existe al menos una coincidencia; si no, \texttt{false}.
      \item \textbf{Ultra-eficiente:} Genera \texttt{SELECT COUNT(id) > 0}, sin traer entidades pesadas a memoria.
    \end{itemize}
  }
  {Rendimiento}
  {Ideal para validaciones previas de disponibilidad de correos o c\'odigos antes de registrar.}

\begin{frame}[fragile]{Ejemplo 2: Verificaci\'on de Existencia con boolean}
\begin{lstlisting}[style=SpringCodeStyle]
// 1. Declaracion en UsuarioRepository
@Repository
public interface UsuarioRepository extends JpaRepository<Usuario, Long> {
    boolean existsByCorreoInstitucional(String correo);
}

// 2. Invocacion en UsuarioService para validaciones livianas
@Service
@RequiredArgsConstructor
public class UsuarioService {
    private final UsuarioRepository usuarioRepository;

    public boolean existeCorreo(String correo) {
        return usuarioRepository.existsByCorreoInstitucional(correo);
    }
}
\end{lstlisting}
\end{frame}

% ------------------------------------------------------------------------------
% PASO 3: BÚSQUEDA DE COLECCIONES (findBy + List$<$>$)
% ------------------------------------------------------------------------------
\contentSlideBulletCard{Colecciones de Entidades: findBy... con List$<$T$>$}
  {
    \begin{itemize}
      \item \textbf{Prop\'osito:} Recuperar todas las entidades que coincidan con un criterio com\'un.
      \item \textbf{Retorno List$<$T$>$:} Contiene todas las coincidencias ordenadas o en bruto.
      \item \textbf{Colecci\'on Vac\'ia:} Si no hay registros coincidentes, Spring retorna una lista vac\'ia (nunca \texttt{null}).
    \end{itemize}
  }
  {M\'ultiples Filas}
  {Ideal para filtrar entidades que comparten una categor\'ia, departamento o estado com\'un.}

\begin{frame}[fragile]{Ejemplo 3: B\'usqueda de Colecciones con List$<$T$>$}
\begin{lstlisting}[style=SpringCodeStyle]
// 1. Declaracion en CursoRepository
@Repository
public interface CursoRepository extends JpaRepository<Curso, Long> {
    List<Curso> findByDepartamento(String departamento);
}

// 2. Invocacion en CursoService
@Service
@RequiredArgsConstructor
public class CursoService {
    private final CursoRepository cursoRepository;

    public List<Curso> listarPorDepartamento(String depto) {
        return cursoRepository.findByDepartamento(depto);
    }
}
\end{lstlisting}
\end{frame}

% ------------------------------------------------------------------------------
% PASO 4: CONTEO AGREGADO (countBy + long)
% ------------------------------------------------------------------------------
\contentSlideBulletCard{Conteo Num\'erico: countBy... con long}
  {
    \begin{itemize}
      \item \textbf{Prop\'osito:} Conocer el total exacto de registros que cumplen una condici\'on.
      \item \textbf{Retorno long:} Devuelve el n\'umero escalar directamente desde el motor SQL.
      \item \textbf{Cero Sobrecarga:} Evita el antipatr\'on de hacer \texttt{findAll().size()} que colapsa la memoria JVM.
    \end{itemize}
  }
  {M\'etricas}
  {Genera directamente \texttt{SELECT COUNT(*)} en base de datos de manera at\'omica y veloz.}

\begin{frame}[fragile]{Ejemplo 4: Conteo Num\'erico con countBy y long}
\begin{lstlisting}[style=SpringCodeStyle]
// 1. Declaracion en CursoRepository
@Repository
public interface CursoRepository extends JpaRepository<Curso, Long> {
    long countByDepartamento(String departamento);
}

// 2. Invocacion en CursoService
@Service
@RequiredArgsConstructor
public class CursoService {
    private final CursoRepository cursoRepository;

    public long contarCursosPorDepartamento(String depto) {
        return cursoRepository.countByDepartamento(depto);
    }
}
\end{lstlisting}
\end{frame}

% ------------------------------------------------------------------------------
% PASO 5: OPERADORES LÓGICOS Y BANDERAS (And, Or, ActiveTrue)
% ------------------------------------------------------------------------------
\contentSlideBulletCard{Filtros Combinados: Operadores And / Or y Banderas}
  {
    \begin{itemize}
      \item \textbf{Operador And:} Exige que se cumplan m\'ultiples condiciones simult\'aneas en el \texttt{WHERE}.
      \item \textbf{Operador Or:} Permite coincidencia de cualquiera de las condiciones especificadas.
      \item \textbf{Banderas Booleanas:} Sufijos directos como \texttt{ActiveTrue} o \texttt{ActiveFalse}.
    \end{itemize}
  }
  {L\'ogica Booleana}
  {Permite componer filtros complejos concatenando nombres de atributos de forma declarativa.}

\begin{frame}[fragile]{Ejemplo 5: Filtro Combinado con And y ActiveTrue}
\begin{lstlisting}[style=SpringCodeStyle]
// 1. Declaracion en ProfesorRepository
@Repository
public interface ProfesorRepository extends JpaRepository<Profesor, Long> {
    List<Profesor> findByDepartamentoAndActiveTrue(String depto);
}

// 2. Invocacion en ProfesorService
@Service
@RequiredArgsConstructor
public class ProfesorService {
    private final ProfesorRepository profesorRepository;

    public List<Profesor> listarProfesoresActivos(String depto) {
        return profesorRepository.findByDepartamentoAndActiveTrue(depto);
    }
}
\end{lstlisting}
\end{frame}

% ------------------------------------------------------------------------------
% PASO 6: RANGOS Y DESIGUALDADES (Between, GreaterThanEqual)
% ------------------------------------------------------------------------------
\contentSlideBulletCard{Rangos y Comparaciones Num\'ericas}
  {
    \begin{itemize}
      \item \textbf{Between:} Rango cerrado inclusivo (\texttt{campo BETWEEN min AND max}).
      \item \textbf{GreaterThan / GreaterThanEqual:} Valores mayores o mayores e iguales que un umbral.
      \item \textbf{LessThan / LessThanEqual:} Valores menores o menores e iguales que un l\'imite.
    \end{itemize}
  }
  {Filtros de Escala}
  {Ideal para consultar rangos de cr\'editos acad\'emicos, edades, fechas o presupuestos.}

\begin{frame}[fragile]{Ejemplo 6: Rangos de Cr\'editos con Between}
\begin{lstlisting}[style=SpringCodeStyle]
// 1. Declaracion en CursoRepository
@Repository
public interface CursoRepository extends JpaRepository<Curso, Long> {
    List<Curso> findByCreditosBetween(int minCreditos, int maxCreditos);
    List<Curso> findByCreditosGreaterThanEqual(int creditosMinimos);
}

// 2. Invocacion en CursoService
@Service
@RequiredArgsConstructor
public class CursoService {
    private final CursoRepository cursoRepository;

    public List<Curso> listarPorRangoCreditos(int min, int max) {
        return cursoRepository.findByCreditosBetween(min, max);
    }
}
\end{lstlisting}
\end{frame}

% ------------------------------------------------------------------------------
% PASO 7: BÚSQUEDA DE TEXTO Y ORDENAMIENTO (Containing, OrderBy)
% ------------------------------------------------------------------------------
\contentSlideBulletCard{B\'usqueda Parcial de Texto y Ordenamiento}
  {
    \begin{itemize}
      \item \textbf{Containing:} Traduce a \texttt{LIKE \%texto\%} para b\'usqueda de fragmentos o palabras clave.
      \item \textbf{IgnoreCase:} Aplica \texttt{UPPER()} en SQL para comparaciones insensibles a may\'usculas.
      \item \textbf{OrderBy...Asc / Desc:} Aplica cl\'ausula \texttt{ORDER BY} directamente en la base de datos.
    \end{itemize}
  }
  {Buscadores Web}
  {Los par\'ametros se procesan v\'ia PreparedStatement, previniendo vulnerabilidades de inyecci\'on SQL.}

\begin{frame}[fragile]{Ejemplo 7: B\'usqueda de Texto y Ordenamiento}
\begin{lstlisting}[style=SpringCodeStyle]
// 1. Declaraciones en repositorios
public interface CursoRepository extends JpaRepository<Curso, Long> {
    List<Curso> findByNombreContainingIgnoreCase(String fragmento);
}

public interface ProfesorRepository extends JpaRepository<Profesor, Long> {
    List<Profesor> findByEspecialidadIgnoreCaseOrderByApellidoAsc(String esp);
}

// 2. Invocacion en Services correspondientes
public List<Curso> buscarCursosPorTexto(String texto) {
    return cursoRepository.findByNombreContainingIgnoreCase(texto);
}

public List<Profesor> listarPorEspecialidadOrdenados(String esp) {
    return profesorRepository.findByEspecialidadIgnoreCaseOrderByApellidoAsc(esp);
}
\end{lstlisting}
\end{frame}

% ------------------------------------------------------------------------------
% PASO 8: NAVEGACIÓN DE RELACIONES (_)
% ------------------------------------------------------------------------------
\contentSlideBulletCard{Navegaci\'on de Relaciones del Modelo (\_)}
  {
    \begin{itemize}
      \item \textbf{Guion Bajo (\_):} Permite descender a trav\'es de relaciones JPA sin ambig\"uedad.
      \item \textbf{Relaciones @ManyToOne:} \texttt{findByProfesor\_Id(Long id)} o por atributos anidados.
      \item \textbf{Relaciones @ManyToMany:} \texttt{findByRoles\_Nombre(String rol)}.
      \item \textbf{Tablas Intermedias:} \texttt{existsById\_EstudianteId(Long estudianteId)}.
    \end{itemize}
  }
  {JOIN Autom\'atico}
  {Spring Data detecta las relaciones mapeadas y genera autom\'aticamente el \texttt{INNER JOIN} relacional.}

\begin{frame}[fragile]{Ejemplo 8: Navegaci\'on de Relaciones con Guion Bajo (\_)}
\begin{lstlisting}[style=SpringCodeStyle]
// 1. Repositorios navegando relaciones ManyToOne y ManyToMany
public interface CursoRepository extends JpaRepository<Curso, Long> {
    List<Curso> findByProfesor_Id(Long profesorId);
    List<Curso> findByProfesor_Departamento(String depto);
}

public interface UsuarioRepository extends JpaRepository<Usuario, Long> {
    List<Usuario> findByRoles_Nombre(String nombreRol);
}

// 2. Invocacion en UsuarioService
public List<Usuario> listarUsuariosPorRol(String rol) {
    return usuarioRepository.findByRoles_Nombre(rol);
}
\end{lstlisting}
\end{frame}

% ==============================================================================
% SECCIÓN 3: CONSULTAS AVANZADAS CON @QUERY (JPQL Y NATIVE QUERIES)
% ==============================================================================

\sectionSlideStripe{03}{Consultas con @Query}{JPQL, Native SQL y Modificaciones Masivas}

\contentSlideBulletCard{?`Cu\'ando los Query Methods no son Suficientes?}
  {
    \begin{itemize}
      \item \textbf{M\'etodos Kilom\'etricos:} Nombres excesivamente largos y dif\'iciles de mantener.
      \item \textbf{Agregaciones Complejas:} Cl\'ausulas \texttt{COUNT}, \texttt{AVG}, \texttt{GROUP BY}, \texttt{HAVING}.
      \item \textbf{JOINs Expl\'icitos:} Consultas que cruzan tablas intermedias asociativas.
    \end{itemize}
  }
  {Control Total}
  {La anotaci\'on \texttt{@Query} permite definir exactamente la consulta requerida.}

\contentSlideTwoCols{JPQL vs Native Query}
  {JPQL (Java Persistence Query Language)}
  {
    \begin{itemize}
      \item Basado en clases y atributos Java.
      \item 100\% portable entre motores SQL.
      \item Valida sintaxis al arrancar la app.
      \item Mapea entidades directamente.
    \end{itemize}
  }
  {Native Query (SQL del Motor)}
  {
    \begin{itemize}
      \item Basado en tablas y columnas f\'isicas.
      \item Acoplado al motor (PostgreSQL/H2).
      \item Permite funciones espec\'ificas y CTEs.
      \item Se valida al momento de ejecutarse.
    \end{itemize}
  }

\contentSlideBulletCard{Par\'ametros Nombrados y Operaciones @Modifying}
  {
    \begin{itemize}
      \item \textbf{Par\'ametros Nombrados:} Se vinculan usando \texttt{:nombre} y la anotaci\'on \texttt{@Param("nombre")}.
      \item \textbf{Sentencias UPDATE/DELETE:} Deben marcarse obligatoriamente con \texttt{@Modifying}.
      \item \textbf{Retorno de Filas:} Retornan un valor primitivo \texttt{int} con el total de filas modificadas.
      \item \textbf{Contexto Transaccional:} Operaciones de modificaci\'on requieren transacci\'on activa.
    \end{itemize}
  }
  {Actualizaci\'on Masiva}
  {\texttt{@Modifying} actualiza directamente en la base de datos sin cargar entidades a memoria.}

\begin{frame}[fragile]{Ejemplo 9: Consulta con JPQL y Native SQL}
\begin{lstlisting}[style=SpringCodeStyle]
@Repository
public interface EstudianteRepository extends JpaRepository<Estudiante, Long> {

    // 1. JPQL: Orientado al modelo de entidades (EstudianteCurso)
    @Query("SELECT ec.estudiante FROM EstudianteCurso ec " +
           "WHERE ec.curso.id = :cursoId AND ec.estudiante.active = true " +
           "ORDER BY ec.estudiante.apellido ASC")
    List<Estudiante> buscarEstudiantesPorCursoJPQL(@Param("cursoId") Long cursoId);

    // 2. Native Query: SQL directo con INNER JOIN sobre tablas fisicas
    @Query(value = "SELECT e.* FROM estudiante e " +
                   "INNER JOIN estudiante_curso ec ON e.id = ec.estudiante_id " +
                   "WHERE ec.curso_id = :cursoId AND e.active = true " +
                   "ORDER BY e.apellido ASC", nativeQuery = true)
    List<Estudiante> buscarEstudiantesPorCursoNativo(@Param("cursoId") Long cursoId);
}
\end{lstlisting}
\end{frame}

% ==============================================================================
% SECCIÓN 4: ESTUDIO AUTÓNOMO (TRANSACCIONES, ACID Y PAGINACIÓN)
% ==============================================================================

\sectionSlideDark{04}{Estudio Aut\'onomo}{Transacciones ACID, Rollback y Paginaci\'on}

\contentSlideQuoteStat{Propiedades ACID}{ACID}{Garant\'ia de Integridad}{“O se ejecutan todas las operaciones con \'exito, o ninguna deja huella en la base de datos.”}

\contentSlideBulletCard{La Anotaci\'on @Transactional}
  {
    \begin{itemize}
      \item \textbf{Proxy Din\'amico (AOP):} Spring intercepta el m\'etodo y gestiona la conexi\'on JDBC.
      \item \textbf{Delimitaci\'on At\'omica:} Abre transacci\'on al entrar y confirma con \texttt{commit}.
      \item \textbf{Gesti\'on de Fallos:} Si ocurre una excepci\'on no controlada, revierte con \texttt{rollback}.
    \end{itemize}
  }
  {Consistencia}
  {Evita estados intermedios inconsistentes entre m\'ultiples tablas ante fallos imprevistos.}

\contentSlideTwoCols{Pol\'iticas de Rollback y Rendimiento}
  {Comportamiento de Rollback}
  {
    \begin{itemize}
      \item \textbf{Por defecto:} Rollback en \texttt{RuntimeException} y \texttt{Error}.
      \item \textbf{Checked Exceptions:} Por defecto NO provocan rollback.
      \item \textbf{Regla robusta:} Usar \texttt{@Transactional(rollbackFor = Exception.class)}.
    \end{itemize}
  }
  {Optimizaci\'on readOnly}
  {
    \begin{itemize}
      \item Con \texttt{readOnly = true}:
      \item Hibernate apaga el \textit{Dirty Checking}.
      \item Ahorra CPU y memoria Heap.
      \item Permite enrutar a r\'eplicas de lectura.
    \end{itemize}
  }

\begin{frame}[fragile]{Ejemplo Aut\'onomo: Transacci\'on y Rollback}
\begin{lstlisting}[style=SpringCodeStyle]
@Service
@RequiredArgsConstructor
public class MatriculaService {
    private final EstudianteRepository estudianteRepo;
    private final CursoRepository cursoRepo;
    private final EstudianteCursoRepository estudianteCursoRepo;

    @Transactional(rollbackFor = Exception.class)
    public EstudianteCurso matricular(Long estId, Long cursoId) {
        var est = estudianteRepo.findById(estId)
            .orElseThrow(() -> new IllegalArgumentException("Estudiante no existe"));
        var cur = cursoRepo.findById(cursoId)
            .orElseThrow(() -> new IllegalArgumentException("Curso no existe"));

        var mat = estudianteCursoRepo.save(new EstudianteCurso(est, cur));

        // Regla tardia: si excede creditos, lanza excepcion y dispara ROLLBACK
        if (cur.getCreditos() > 4) {
            throw new RuntimeException("Creditos excedidos -> ROLLBACK de la matricula");
        }
        return mat;
    }
}
\end{lstlisting}
\end{frame}

\contentSlideBulletCard{Paginaci\'on: ?`Por qu\'e Evitar findAll()?}
  {
    \begin{itemize}
      \item \textbf{Saturaci\'on de Memoria:} Cargar miles de filas satura el Heap (OutOfMemory).
      \item \textbf{Latencia de Red:} Transferir payloads JSON gigantescos bloquea al cliente.
      \item \textbf{Carga en BD:} Bloqueos prolongados de lectura y alto I/O en disco.
    \end{itemize}
  }
  {Escalabilidad}
  {La paginaci\'on fragmenta grandes vol\'umenes de datos entregando bloques controlados bajo demanda.}

\contentSlideTwoCols{Page$<$T$>$ vs Slice$<$T$>$}
  {Page$<$T$>$ (Paginado Completo)}
  {
    \begin{itemize}
      \item Ejecuta \texttt{SELECT ... LIMIT} y \texttt{COUNT(*)}.
      \item Conoce el total de p\'aginas y registros.
      \item Ideal para interfaces web con barra de p\'aginas (1, 2, 3...).
    \end{itemize}
  }
  {Slice$<$T$>$ (Scroll Infinito)}
  {
    \begin{itemize}
      \item Ejecuta \texttt{SELECT ... LIMIT + 1}.
      \item S\'olo sabe si existe p\'agina siguiente.
      \item Ahorra el costo de contar todo con \texttt{COUNT(*)}.
      \item Ideal para aplicaciones m\'oviles y feeds infinitos.
    \end{itemize}
  }

\begin{frame}[fragile]{Ejemplo Aut\'onomo: Paginaci\'on y Ordenamiento}
\begin{lstlisting}[style=SpringCodeStyle]
// Repositorio: Spring Data inyecta LIMIT y OFFSET segun el Pageable
public interface CursoRepository extends JpaRepository<Curso, Long> {
    Page<Curso> findByDepartamento(String depto, Pageable pageable);
}

// Servicio: construye el PageRequest con paginacion y orden
@Transactional(readOnly = true)
public Page<Curso> listarPaginados(String depto, int page, int size, String sort) {
    Pageable pageable = PageRequest.of(page, size, Sort.by(sort).ascending());
    return cursoRepository.findByDepartamento(depto, pageable);
}

// Controlador REST: expone los parametros de consulta al cliente
@GetMapping("/paginados")
public ResponseEntity<Page<Curso>> getPaginados(
        @RequestParam String depto, @RequestParam(defaultValue = "0") int page,
        @RequestParam(defaultValue = "10") int size, @RequestParam(defaultValue = "nombre") String sort) {
    return ResponseEntity.ok(cursoService.listarPaginados(depto, page, size, sort));
}
\end{lstlisting}
\end{frame}

\contentSlideQuoteStat{Conclusi\'on de la Sesi\'on}{100\%}{Arquitectura Profesional}{“El dise\~no riguroso de Query Methods y @Query en la persistencia, coordinado por Servicios limpios, garantiza un backend expresivo, desacoplado y escalable.”}

\end{document}
